\section{Proof of Proposition~\ref{prop:identification}}
\label{app:identification}
We give the proof for a fixed latent representation. The observed-design statement follows by the same argument, as explained at the end. All conditional expectations in the first three steps are given $(\bm{D}_1,\bm{D}_2)$. Write $G_i(t)=G(t\mid \bm{D}_i)$, $\bm{w}=\bm{D}_1-\bm{D}_2$, and $\bm{v}=\bm{\vartheta}-\bm{\vartheta}_0$.

\paragraph{Event-weighted differences.}
On $\{\Delta_i=1\}$, $R_i=\varepsilon_i$, so the conditional event subdensity is $f_\varepsilon(t)G_i(t)$. Moreover,
\[
 \Pr(R_j\ge t\mid \bm{D}_j)=S_\varepsilon(t)G_j(t).
\]
It follows that
\begin{equation}
\begin{aligned}
A_{12}&:=E\{\Delta_1\mathbf1(R_2\ge R_1)\mid\bm D_1,\bm D_2\}\\
&=\int f_\varepsilon(t)S_\varepsilon(t)G_1(t)G_2(t)\,\mathrm{d}t=A_{21}.
\end{aligned}
\label{eq:app-centering}
\end{equation}
The finite measure $\Pr(\Delta_1=1,R_2-R_1\in\mathrm{d}r\mid \bm{D}_1,\bm{D}_2)$ has a density, say $q_{12}(r)$. With $F_{R_2\mid \bm{D}_2}$ denoting the conditional distribution of $R_2$,
\begin{equation}
 q_{12}(r)=\int f_\varepsilon(z-r)G_1(z-r)\,\mathrm{d}F_{R_2\mid \bm{D}_2}(z)
 \le L_\varepsilon\qquad\text{a.e.}
 \label{eq:app-density}
\end{equation}
The Stieltjes integral allows censoring-induced atoms in $R_2$. Restricting it to the event submeasure of patient 2 gives, for $|r|\le A$,
\begin{equation}
 \begin{aligned}
 q_{12}(r)
 &\ge \int f_\varepsilon(z-r)G_1(z-r)
              f_\varepsilon(z)G_2(z)\,\mathrm{d}z\\
 &\ge (2A-|r|)m_\varepsilon^2g_0^2
 \ge c_{\mathrm{loc}},
 \end{aligned}
 \label{eq:app-density-lower}
\end{equation}
because $I\cap(I+r)$ has length $2A-|r|$. These bounds hold almost everywhere, which suffices for the integrations below.

\paragraph{A finite conditional contrast.}
Set $t=\bm{w}^\top \bm{v}$. At the candidate parameter,
\[
 e_2^o(\bm{\vartheta})-e_1^o(\bm{\vartheta})=R_2-R_1+t.
\]
Define
\[
 k_{12}(t)=E\!\left[\Delta_1
       \{(R_2-R_1+t)_+-(R_2-R_1)_+\}\mid \bm{D}_1,\bm{D}_2\right].
\]
This is finite, even without an error moment, because its absolute value is at most $|t|$. The absolute continuity of the event-weighted difference measure gives $k'_{12}(0)=A_{12}$ and $k''_{12}(t)=q_{12}(-t)$ almost everywhere. Integrating this derivative, for $t\ge0$,
\[
 k_{12}(t)-A_{12}t=\int_0^t(t-z)q_{12}(-z)\,\mathrm{d}z.
\]
For $t<0$, the corresponding expression is
\[
 k_{12}(t)-A_{12}t=\int_t^0(z-t)q_{12}(-z)\,\mathrm{d}z.
\]
Nonnegativity of the density and~\eqref{eq:app-density-lower} imply in either case
\begin{equation}
 k_{12}(t)-A_{12}t\ge\frac{c_{\mathrm{loc}}}{2}\rho_A(t).
 \label{eq:app-hinge-margin}
\end{equation}
For example, when $t>A$ the retained integral is
$c_{\mathrm{loc}}\int_0^A(t-z)\,\mathrm{d}z
 =c_{\mathrm{loc}}(At-A^2/2)$; the case $0\le t\le A$ gives $c_{\mathrm{loc}}t^2/2$. The negative cases follow analogously.

\paragraph{Cancellation and strict separation.}
Since the two patients are identically distributed, interchanging their indices preserves their joint law, changes $\bm{w}$ to $-\bm{w}$, and leaves $A_{12}$ unchanged by~\eqref{eq:app-centering}. The integrability of $\bm{w}^\top \bm{v}$ therefore gives
\[
 E(A_{12}\bm{w}^\top \bm{v})=0.
\]
Taking expectations in~\eqref{eq:app-hinge-margin} proves~\eqref{eq:population-gap}. When $\bm{\vartheta},\bm{\vartheta}_0\in\Theta_s$, their difference obeys $\|\bm{v}_\beta\|_0\le2s$. Assumption~\ref{ass:separation} yields
\[
 E\{(\bm{w}^\top \bm{v})^2\}=2\operatorname{Var}(\bm{D}_i^\top \bm{v})
 \ge2\kappa_{\mathrm{pop}}\|\bm{v}\|_2^2.
\]
For nonzero $\bm{v}$, $\Pr(\bm{w}^\top \bm{v}\ne0)>0$. Because $\rho_A$ vanishes only at zero, the contrast is strictly positive. Also $\rho_A(t)\le 2A|t|$, so the expectation in~\eqref{eq:population-gap} is finite. This proves the unique minimum at $\bm{\vartheta}_0$ on $\Theta_s$ without requiring a finite expected uncentered loss.

\paragraph{The ordinary observed-data target.}
For $\mathcal M_0$, set $\bm{V}_i=(\bm{Z}_i^\top,\bm{X}_i^\top)^\top$ and $\bm{b}=(\bm{\theta}^\top,\bm{\beta}^\top)^\top$. Replace $\bm{w}^\top \bm{v}$ by $(V_1-V_2)^\top(\bm{b}-\bm{b}_0)$ in the preceding calculation. The event-weighted centering and density bounds can still be established conditional on $(\bm{D}_1,\bm{D}_2)$ and then averaged. A positive variance for every nonzero sparse difference in this observed predictor gives strict separation.

Every contrast increment can now be computed as an expectation of a function of the observed pair $(O_1,O_2)$. To avoid an unknown reference parameter in the identification argument, one can fix any deterministic $\bm{b}_\dagger$ and subtract the loss at $\bm{b}_\dagger$ instead. Every such contrast is finite and differs from the one centered at $\bm{b}_0$ by a constant independent of $\bm{b}$. Thus its unique minimizer is a functional of the observed-data law. The corresponding $\mathcal M_1$ result is relative to the specified latent representation.
\hfill$\square$
